\section{Related Work}

\subsection{Local Inference Engines and Single-Model Optimization}

The most relevant adjacent line of work focuses on improving local execution efficiency for a \emph{single} model on constrained hardware.

\textbf{vLLM}\cite{kwon2023vllm} introduces \textit{PagedAttention} to manage KV cache efficiently and improve throughput in high-concurrency serving. However, its primary objective is high-throughput serving for one model, rather than service-level arbitration among heterogeneous tasks sharing a single constrained GPU.

\textbf{FlexGen}\cite{sheng2023flexgen} uses aggressive tensor offloading to CPU and disk to run large models on limited hardware. While effective for capacity-oriented or offline settings, its I/O-heavy design is not aimed at preserving interactive continuity during rapid task switching.

\textbf{Llama.cpp}\cite{llamacpp} enables practical LLM execution on consumer CPUs/GPUs through quantization and mmap-based loading. It is highly relevant as the local inference substrate in our setting, but it is not itself a scheduler for multiple heterogeneous services. Accordingly, our evaluation uses \texttt{llama.cpp + manual switching} as a practical baseline, rather than treating \texttt{llama.cpp} alone as a scheduling system.

Other optimization efforts, including \textbf{PowerInfer}\cite{song2023powerinfer}, \textbf{Deja Vu}\cite{liu2023dejavu}, and \textbf{MLC LLM}\cite{wu2023mlc}, improve single-model execution via sparsity or compilation. These systems are complementary to our work: they optimize the internal execution path of one model, whereas EdgeSched arbitrates the lifecycle and resource ownership of multiple workload classes at the system level.

More recent serving systems move closer to the scheduling boundary that motivates our work. \textbf{Splitwise}\cite{patel2024splitwise} and \textbf{P/D-Serve}\cite{jin2024pdserve} separate prefill and decode phases to improve utilization in high-throughput LLM serving. \textbf{Inference without Interference}\cite{hu2024interference} studies mixed downstream workloads and argues for disaggregating inference paths to reduce cross-workload interference. These systems are conceptually closer to EdgeSched than pure single-model optimizers because they recognize that different phases or request classes may require different control strategies. However, their design center remains server-style LLM serving with phase-aware resource allocation, rather than single-device arbitration between a resident service and an unrelated GPU-exclusive application.

This distinction matters for evaluation: if one wanted stronger external comparators than manual switching or OS-managed memory, systems in the Splitwise/P-D-Serve/Inference-without-Interference family would be the most credible conceptual references. At the same time, they are not drop-in baselines for our deployment because they assume LLM-serving stacks, batching-oriented request flows, or distributed/server-class execution environments that do not match our LLM-plus-image-generation setting.

\subsection{Application-Level Memory Arbitration}

Traditional OS scheduling provides mature mechanisms for CPU processes, but it offers limited control over GPU memory ownership at the application level. Under dedicated-VRAM constraints, relying purely on OS-level paging, swap-like behavior, or unified-memory style fallback can lead to severe thrashing or timeout when multiple VRAM-intensive workloads compete. Our evaluation confirms this effect through the LRU/unified-memory baseline.

EdgeSched addresses a different question from single-model inference engines: how to arbitrate between a \emph{resident, migratable, latency-sensitive} workload and a \emph{bursty, GPU-pinned, queueable} workload on a single consumer device. The key contribution is therefore not a new low-level primitive, but a control consequence of service-role asymmetry: under the constraints studied here, the scheduler should release VRAM quickly on the downlink path while deferring visible interruption on the uplink path.

This is also where EdgeSched differs from a mere engineering bundle. Priority queues, offloading, and crash recovery each exist elsewhere; the paper's claim is that the studied local-runtime setting imposes different control objectives on the two migration directions, which in turn motivates separating urgent downlink from continuity-oriented restoration.

\subsection{General-Purpose Orchestration Systems}

Frameworks such as \textbf{Ray}\cite{moritz2018ray} and \textbf{Kubernetes}\cite{verma2015borg} are designed for distributed orchestration across clusters or containers. They provide general-purpose scheduling abstractions, fault handling, and elastic resource management, but their design center is multi-node coordination rather than sub-second resource arbitration on a single consumer device. Even edge extensions such as \textbf{KubeEdge}\cite{xiong2019kubeedge} inherit much of the architectural overhead of cloud-native orchestration.

We mention these systems only to clarify scope: EdgeSched does not target distributed placement or cluster-level orchestration. Instead, it targets single-node local AI runtimes with limited dedicated VRAM. Therefore, Ray/Kubernetes-style systems are relevant primarily as \emph{positioning references}, not as direct experimental baselines in our deployment setting.

\subsection{Comparison Summary}

Table~\ref{tab:comparison} summarizes the role of representative prior systems relative to EdgeSched.

\begin{table}[htbp]
\caption{Comparison of Representative Systems}
\label{tab:comparison}
\centering
\scriptsize
\resizebox{\columnwidth}{!}{%
\begin{tabular}{@{}lccccc@{}}
\toprule
System & Primary & Deploy. & Sched. & Multi-Workload & Overhead \\
 & Focus & Scope & Level & Arbitration & Profile \\
\midrule
Ray & Dist. orchestration & Cluster & Coarse & Yes & High \\
Kubernetes & Container orchestration & Cluster & Coarse & Yes & High \\
vLLM & Single-model throughput & Server & Intra-model & No & Medium \\
FlexGen & Capacity via offloading & Single model & Intra-model & No & Medium \\
Splitwise / P-D-Serve & Phase-disagg. serving & Server & Inter-phase & Partial & Medium \\
Inference w/o Interference & Mixed LLM workloads & Server & Inter-service & Partial & Medium \\
Llama.cpp & Local inference & Consumer & None & No & Low \\
\textbf{EdgeSched} & Local AI scheduling & Consumer & Inter-service & Yes & \textbf{Low} \\
\bottomrule
\end{tabular}%
}
\end{table}
